\chapter{Kluczowe mechanizmy i~algorytmy}

\section{Motywy kolorystyczne i~kontrast WCAG}

\subsection{Pięć tokenów bazowych}
Cały wygląd witryny wynika z~pięciu zmiennych CSS zdefiniowanych w~\code{:root} pliku
\code{static/css/main.css}: \code{--color-primary}, \code{--color-secondary},
\code{--color-accent}, \code{--color-surface}, \code{--color-surface-variant}. Wszystkie
pozostałe odcienie (najechanie, ramki, tła kart, teksty pomocnicze) są wyprowadzane
funkcjami \code{color-mix()} lub składnią względną \code{hsl(from …)}, np.:

\begin{lstlisting}
--color-primary-light: color-mix(in srgb, var(--color-primary) 22%, var(--color-surface));
--color-primary-hover: color-mix(in srgb, var(--color-primary) 85%, black);
--color-on-surface-muted: color-mix(in srgb, var(--color-on-surface) 74%, var(--color-surface));
\end{lstlisting}

W~arkuszu nie ma innych literałów kolorów (wyjątek: logotypy portali OAuth). W~stanie domyślnym szablon bazowy nie nadpisuje \code{:root}, więc
ręczna zmiana tokenu w~pliku CSS zmienia tonację całej witryny. Gdy użytkownik wybierze
motyw, \code{base.html} wstawia blok \code{<style>} z~pięcioma kolorami motywu i~wyliczonymi
kolorami tekstu.

\subsection{Wyliczanie koloru tekstu}
Kolor tekstu na każdym z~teł (\code{--color-on-*}) i~tryb jasny/ciemny
(\code{data-theme-mode}) wylicza moduł \code{apps/core/colors.py} zgodnie z~definicją WCAG
2.1. Dla kanału $c \in \{R,G,B\}$ zapisanego w~zakresie 0--255:
\[
c_{\text{lin}} = \begin{cases} \dfrac{c/255}{12{,}92} & \text{gdy } c/255 \le 0{,}04045\\[8pt]
\left(\dfrac{c/255 + 0{,}055}{1{,}055}\right)^{2{,}4} & \text{w~przeciwnym razie}\end{cases}
\qquad L = 0{,}2126\,R_{\text{lin}} + 0{,}7152\,G_{\text{lin}} + 0{,}0722\,B_{\text{lin}}
\]
\[
\text{kontrast}(C_1, C_2) = \frac{L_{\max} + 0{,}05}{L_{\min} + 0{,}05} \in [1, 21]
\]
Funkcja \code{readable\_text\_color(tło)} wybiera biały albo prawie czarny
(\code{\#111827}) tekst -- ten, który daje wyższy kontrast z~tłem. Motyw jest oznaczany
jako ciemny, gdy tekst na kolorze \code{surface} ma być jasny. Panel administracyjny
(django-unfold) otrzymuje skalę odcieni wyliczoną funkcją \code{unfold\_palette()}, która
szuka najsłabszego przesunięcia koloru spełniającego kontrast 4,5:1 na wszystkich tłach.
Administrator wpisuje więc tylko pięć kolorów -- czytelność gwarantuje system.

\subsection{Kontrola dostępności motywu}\label{sec:theme-check}
Formularz motywu w~panelu administratora (\code{ColorThemeForm.clean}) odrzuca zestaw kolorów,
dla którego \code{theme\_accessibility\_problems()} zgłasza problem: tekst na tle i~na tle
kontenerów poniżej 4,5:1 (lub tła wymagające różnych kolorów tekstu), tekst pomocniczy (76\%
tekstu + tło) poniżej 4,5:1, napis na przycisku w~kolorze głównym, dodatkowym lub akcentowym
poniżej 4,5:1 oraz kolor marki użyty jako tekst poniżej 4,5:1. Komunikaty po polsku podają
zmierzony kontrast i~wskazówkę, co zmienić.

\subsection{Rozstrzyganie aktywnego motywu}
Procesor kontekstu \code{theme\_colors} wybiera motyw w~kolejności: motyw zapisany
w~profilu zalogowanego użytkownika $\rightarrow$ motyw z~sesji gościa $\rightarrow$
motyw oznaczony jako domyślny $\rightarrow$ wartości z~arkusza CSS.

\section{Uprawnienia moderatorów w~drzewie kategorii}

Nauczyciel moderuje kategorię, jeśli jest przypisany do niej \emph{lub do dowolnego jej
przodka}. Dzięki MPTT sprawdzenie sprowadza się do jednego zapytania: kategoria
\code{c} jest w~gałęzi kategorii \code{m}, gdy \code{c.tree\_id = m.tree\_id} oraz
\code{m.lft $\le$ c.lft} i~\code{c.rght $\le$ m.rght}. Na tej samej zasadzie działa
filtr katalogu ,,kategoria z~podkategoriami'' i~liczniki kursów w~drzewie. Lista osób
powiadamianych o~zgłoszeniu (\code{course\_reviewers}) to moderatorzy wszystkich przodków
kategorii kursu, a~gdy ich brak -- administratorzy.

\section{Reputacja i~rangi}

Każdy głos \code{CommentVote} ma wartość $+1$ lub $-1$. Sygnały \code{post\_save}
i~\code{post\_delete} przeliczają reputację autora komentarza jako sumę głosów na
wszystkie jego komentarze (przeliczenie zamiast inkrementacji odporne jest na zmianę głosu
i~usunięcie komentarza), a~następnie ustalają rangę:

\begin{table}[H]
\centering\small
\caption{Progi rang}
\begin{tabular}{@{}lccccc@{}}
\toprule
\hd{Ranga} & Obserwator & Nowicjusz & Pomocnik & Ekspert & Mistrz \\
\midrule
\hd{Punkty} & $\le -1$ & 0 -- 10 & 11 -- 50 & 51 -- 200 & $> 200$ \\
\bottomrule
\end{tabular}
\end{table}

Ponowne kliknięcie tego samego przycisku cofa głos, kliknięcie przeciwnego zmienia jego
wartość. Nie można głosować na własne ani nieopublikowane komentarze.

\section{Ocenianie testów}

Dla pytania zamkniętego z~liczbą punktów $p$, zbiorem poprawnych odpowiedzi $K$ i~zbiorem
zaznaczonym przez ucznia $Z$ przyznawane jest
\[
\text{punkty} = \begin{cases} p & \text{gdy } Z = K \\ 0 & \text{w~przeciwnym razie,}\end{cases}
\]
czyli w~pytaniu wielokrotnego wyboru nie ma punktów częściowych -- zaznaczenie
dodatkowej błędnej odpowiedzi zeruje wynik pytania. Pytanie otwarte otrzymuje punkty od
oceniającego (0..$p$). Wynik procentowy podejścia to
$100 \cdot \sum \text{punkty} / \sum p$, a~test jest zaliczony, gdy wynik
$\ge$ \code{pass\_percent}. Gdy test zawiera choć jedno pytanie otwarte, wynik jest
ustalany dopiero po ocenie ręcznej (\code{finalize\_attempt}).

\paragraph{Blokady w~programie.} Funkcja \code{build\_outline()} przechodzi program kursu
w~kolejności (sekcje, lekcje, testy po lekcjach, testy końcowe). Po napotkaniu
niezaliczonego testu obowiązkowego wszystkie kolejne elementy otrzymują flagę
\code{locked}. Autor, nauczyciele kategorii i~administrator w~trybie podglądu nie podlegają
blokadom; jeśli jednak sami zapisali się na kurs, domyślnie uczą się jak uczniowie
i~mogą przełączyć tryb.

\paragraph{Postęp i~ukończenie.} Postęp kursu to odsetek ukończonych elementów (lekcji
i~zaliczonych testów). Gdy osiągnie 100\%, \code{on\_course\_completed} oznacza zapis
jako ukończony, aktualizuje statystyki kursu, wysyła powiadomienie i~-- jeśli kurs ma
włączoną opcję -- wystawia dyplom z~unikalnym numerem. Dyplom PDF generuje
\code{reportlab} z~fontami osadzonymi w~repozytorium; adres weryfikacji
\code{/learning/dyplom/<kod>/} jest publiczny.

\section{Koszyk i~realizacja zamówienia}

\begin{itemize}
  \item \textbf{Kumulacja.} \code{add\_to\_cart} szuka pozycji z~tym samym kursem; jeśli
        istnieje, zwiększa \code{quantity} (z~ograniczeniem do 10), w~przeciwnym razie
        tworzy pozycję. Unikalność (\code{cart}, \code{course}) gwarantuje to również na
        poziomie bazy.
  \item \textbf{Koszyk gościa.} Identyfikator koszyka jest trzymany w~sesji. Sygnał
        \code{user\_logged\_in} scala go z~koszykiem konta (\code{merge\_carts} sumuje
        ilości tych samych kursów).
  \item \textbf{Zamówienie.} \code{place\_order} w~jednej transakcji kopiuje ceny
        i~tytuły, nadaje numer \code{ORD-RRRR-NNNNNN} (kolejny w~danym roku), zapisuje
        adres (opcjonalnie także w~profilu), czyści koszyk i~rejestruje wysłanie e-maila po
        zatwierdzeniu transakcji.
  \item \textbf{Realizacja.} Dla każdej pozycji: jeśli kupujący nie ma jeszcze kursu,
        pierwsza sztuka tworzy \code{Enrollment}; pozostałe sztuki tworzą kody
        \code{AccessCode}. Kod można wykorzystać raz, na koncie, które kursu jeszcze nie ma.
  \item \textbf{Bezpieczeństwo płatności.} Webhook sprawdza podpis
        \code{Stripe-Signature}, numer zamówienia z~\code{metadata}, kwotę w~groszach
        i~walutę; zdarzenie dla zamówienia już opłaconego jest ignorowane.
\end{itemize}

\section{Asystent AI}

\paragraph{Prompt.} Wiadomość systemowa (\code{llm.system\_prompt}) określa rolę
asystenta (Bielik -- asystent nauki; odpowiedzi wyłącznie po polsku, zwięźle, zwykle do
200 słów; opieranie się na materiałach kursu i~przyznawanie się do braku informacji zamiast
zmyślania; naprowadzanie zamiast podawania gotowych odpowiedzi do testów; kod w~blokach,
bez HTML) oraz dołącza kontekst strony: treść lekcji oczyszczoną z~HTML (do 6000
znaków), opis i~program kursu albo wynik testu (pytania, odpowiedzi ucznia, punkty,
komentarze nauczyciela; poprawne odpowiedzi tylko wtedy, gdy uczeń mógłby je i~tak
zobaczyć na stronie wyniku). Za nią idzie do 12 ostatnich wiadomości rozmowy, z~pominięciem
odmów.

\paragraph{Blokada w~trakcie testu.} Otwarcie testu do rozwiązania tworzy
\code{QuizSitting}; wysłanie testu go usuwa. Dopóki istnieje rekord młodszy niż
\code{CHATBOT\_QUIZ\_LOCK\_MINUTES} (60 min), każde pytanie -- z~dowolnej strony
i~dowolnego kursu -- otrzymuje stałą odmowę bez kontaktu z~modelem. Uniemożliwia to
proszenie asystenta o~rozwiązanie testu.

\paragraph{Bezpieczeństwo odpowiedzi.} Tekst modelu jest najpierw escapowany, a~dopiero
potem zamieniany z~ograniczonego Markdownu (pogrubienie, kursywa, kod, listy, akapity) na
HTML (\code{formatting.render\_reply}) -- model nie może wstrzyknąć znaczników do strony.

\paragraph{Parametry.} Bielik-Minitron 7B v3.0 Instruct w~kwantyzacji Q4\_K\_M,
temperatura 0,4, okno kontekstu 8192 tokeny, model utrzymywany w~pamięci 30 minut. Na
komputerze z~kartą graficzną pierwszy fragment odpowiedzi pojawia się po kilku sekundach,
a~cała odpowiedź trwała w~testach około 5~s.

\section{Style uczenia się i~rekomendacja stylu}\label{sec:styles-algo}

\paragraph{Wybór wersji lekcji.} Funkcja \code{resolve\_version()} (\code{apps/learning/styles.py})
zwraca wersję dla stylu ucznia, jeśli lekcja ma włączone wersje i~istnieje wersja dla jego
aktywnego stylu; w~przeciwnym razie treść podstawową. Parametr \code{?wersja=<slug>} (lub
\code{podstawowa}) pozwala obejrzeć inną wersję. Treść i~film wersji zastępują treść i~film
lekcji; pliki do pobrania i~testy są wspólne, aby wyniki dało się porównywać między stylami.

\paragraph{Zbieranie danych.} Oznaczenie lekcji jako ukończonej zapisuje w~\code{LessonProgress.style}
wersję, którą uczeń właśnie czytał. Przy wysłaniu testu \code{style\_for\_attempt()} przypisuje
podejściu najczęstszy styl wersji przerobionych przez ucznia w~tym kursie, a~gdy takich nie ma --
styl aktualnie wybrany na koncie.

\paragraph{Rekomendacja.} Dla każdego aktywnego stylu $s$ liczona jest liczba ocenionych podejść
$n_s$ i~ich średni wynik $\bar{x}_s$. Aby pojedynczy test nie przesądzał o~wyniku, średnia jest
ściągana w~stronę ogólnej średniej ucznia $\bar{x}$ (wygładzanie bayesowskie z~wagą $K=2$):
\[
w_s = \frac{n_s\,\bar{x}_s + K\,\bar{x}}{n_s + K}.
\]
Kandydatami są style z~$n_s \ge 2$; wygrywa najwyższe $w_s$ (przy remisie -- odsetek zaliczeń,
potem liczba lekcji w~tej wersji). Jeśli uczeń nie ma stylu, system proponuje zwycięzcę; jeśli
zwycięzca różni się od stylu ucznia o~co najmniej 5 punktów procentowych -- poleca zmianę;
w~przeciwnym razie potwierdza obecny styl. Przy braku danych proponuje styl jeszcze
niesprawdzony, aby zebrać porównanie. Rekomendacja nigdy nie zmienia stylu automatycznie.
Styl ucznia trafia też do promptu asystenta AI, który dopasowuje formę wyjaśnień.

\section{Ankieta stylu i~polecanie kursów}\label{sec:recommend}

\paragraph{Ankieta.} Wynik stylu $s$ to odsetek odpowiedzi wskazujących ten styl:
$a_s = n_s / N$, gdzie $N$ -- liczba pytań (domyślnie 10). Styl dominujący to
$\arg\max_s a_s$; remis rozstrzyga kolejność stylów (pole \code{order}), więc te same odpowiedzi
zawsze dają ten sam wynik. Pytania opisują zachowania (,,Musisz nauczyć się obsługi nowego
programu. Od czego zaczynasz?''), a nie pytają wprost o~styl.

\paragraph{Profil stylów ucznia.} Z~trzech źródeł powstaje wektor udziałów $p_s$ (suma 1):
\[
\tilde p_s = 1{,}0\cdot[s = \text{styl wybrany}] \;+\; 1{,}0\cdot a_s \;+\; 0{,}5\cdot[s = \text{styl polecany z~wyników testów}],
\qquad p_s = \tilde p_s \Big/ \textstyle\sum_t \tilde p_t .
\]
Deklaracja ucznia i~ankieta mają równą wagę, a~wyniki testów -- połowę (to sygnał pośredni, liczony
dopiero po zebraniu danych). Brak któregoś źródła po prostu go pomija.

\paragraph{Dopasowanie kursu.} Dla kursu $k$ i~stylu $s$:
\[
f_{k,s} = 0{,}7\cdot\frac{\text{lekcje kursu z~wersją dla } s}{\text{wszystkie lekcje kursu}} \;+\; 0{,}3\cdot[\text{autor zaznaczył } s],
\qquad \text{match}(k) = \sum_s p_s\, f_{k,s}.
\]
Większą wagę ma to, co kurs faktycznie zawiera (wersje lekcji), mniejszą -- deklaracja autora.
Polecane są opublikowane kursy z~$\text{match}>0$, na które uczeń nie jest zapisany, malejąco
(przy remisie decyduje ocena i~liczba uczestników). Uzasadnienie przy kursie wskazuje styl
o~największym wkładzie, np. ,,4 z~7 lekcji ma wersję dla stylu »Praktyk«''. Wartość liczbowa nie
jest pokazywana uczniowi -- ma znaczenie tylko kolejność.

\paragraph{Ograniczenia metody.} Teoria stylów uczenia się (m.in. VARK) jest krytykowana -- badania
nie potwierdzają jednoznacznie, że nauczanie dopasowane do deklarowanego stylu poprawia wyniki.
Dlatego system nie traktuje deklaracji ani ankiety jako pewnika: porównuje rzeczywiste wyniki testów
po lekcjach w~różnych wersjach (sekcja~\ref{sec:styles-algo}) i~na tej podstawie podpowiada zmianę
stylu, a~polecanie kursów uwzględnia ten sygnał. Wszystkie wagi są heurystykami projektowymi.

\section{Statystyki i~postępy}
Panel uczniów kursu i~statystyki platformy liczą dane na bieżąco z~bazy (agregacje ORM:
\code{Count}, \code{Avg}, \code{Sum}), więc nie wymagają osobnych tabel. Procent ukończenia to
$100 \cdot \text{ukończone zapisy} / \text{wszystkie zapisy}$. Zdenormalizowane liczniki kursu
(\code{students\_count}, \code{completed\_count}) są przeliczane przy zapisie, ukończeniu
i~cofnięciu ukończenia kursu.

\section{Galeria i~slajder}

\paragraph{Skalowanie.} Po walidacji formularza (rozszerzenie, poprawność obrazu w~Pillow,
do 10 MB) \code{GalleryItem.save()} obraca obraz zgodnie z~EXIF, zmniejsza go
proporcjonalnie do mieszczenia się w~1600$\times$1200 px i~zapisuje wymiary przed i~po --
użytkownik dostaje komunikat ,,przeskalowano z~$W\times H$''.

\paragraph{Zmiana kolejności.} Konfigurator wysyła po upuszczeniu karty JSON
\code{\{"items": [id zdjęcia, …]\}} z~nową kolejnością. Serwer sprawdza, czy lista zawiera
dokładnie te same zdjęcia co układ (bez powtórzeń), i~w~transakcji przypisuje
\code{order = 1, 2, 3, …} tylko tym istniejącym rekordom, których pozycja się zmieniła --
żaden rekord nie jest usuwany ani tworzony.

\paragraph{Pętla bez końca.} Slajder (\code{gallery\_slider.js}) nie klonuje kart:
przesunięcie w~prawo animuje przesunięcie ścieżki o~szerokość jednej karty, a~po
zakończeniu animacji przenosi pierwszy element DOM na koniec i~zeruje przesunięcie (dla
kierunku przeciwnego -- odwrotnie). Dzięki temu pętla działa dla dowolnej liczby kart,
także mniejszej niż liczba widocznych. Liczbę widocznych kart (1/2/3) określa zmienna
\code{--per-view} zależna od szerokości ekranu (progi 48 i~64 rem); karty poza widokiem
mają atrybut \code{inert}, więc fokus klawiatury do nich nie trafia.

\section{Powiadomienia}

Wszystkie moduły tworzą powiadomienia jedną funkcją \code{notify(user, kind, title,
message, url)} (lub \code{notify\_many}), która sprawdza rodzaj i~to, czy adres jest
wewnętrzny. Rodzaje: komentarz zaakceptowany / oczekujący / odrzucony, odpowiedź na
komentarz, nowe zgłoszenie kursu, wiadomość i~zmiana statusu zgłoszenia, test do
sprawdzenia, test oceniony, kurs ukończony, zmiana statusu zamówienia, komunikat
systemowy. Otwarcie powiadomienia to POST, który oznacza je jako przeczytane i~przekierowuje
pod zapisany adres.

\section{Mechanizmy bezpieczeństwa -- podsumowanie}

\begin{table}[H]
\centering\small
\caption{Zagrożenia i~zastosowane zabezpieczenia}
\begin{tabularx}{\linewidth}{@{}P{3.4cm}L@{}}
\toprule
\hd{Zagrożenie} & \hd{Zabezpieczenie} \\
\midrule
SQL Injection & wyłącznie ORM Django (zapytania parametryzowane), brak \code{raw()}/\code{extra()} \\
XSS & autoescape szablonów, \code{bleach} z~listą dozwolonych znaczników dla HTML z~CKEditora, escapowanie odpowiedzi modelu \\
CSRF & token CSRF w~każdym formularzu POST, zmiany stanu tylko POST, API wyłącznie z~JWT \\
Podszycie się pod płatność & podpis HMAC webhooka, kontrola kwoty/waluty, weryfikacja sesji przez API \\
Nieuprawniony dostęp & sprawdzanie ról i~własności w~widokach, pliki lekcji i~odpowiedzi przez chronione widoki \\
Złośliwe pliki & walidacja typu i~rozmiaru, ponowne zapisanie obrazów przez Pillow, nazwy plików generowane przez serwer \\
Otwarte przekierowania & przekierowania po logowaniu i~z~powiadomień tylko na adresy wewnętrzne \\
Przejęcie konta przez OAuth & dołączanie portalu do konta dopiero po zalogowaniu na to konto, ten sam e-mail, ważność 15 min \\
Zgadywanie hasła & czasowa blokada logowania: 5 nieudanych prób w~15 min $\rightarrow$ 5 min blokady konta (login i~e-mail liczą się razem); licznik w~Redis \\
Clickjacking & \code{XFrameOptionsMiddleware} \\
\bottomrule
\end{tabularx}
\end{table}
